\section{Methodology}

This section details the technical approach employed to construct and train Named Entity Recognition (NER) models. The methodology focuses on the selection of foundational Transformer architectures, subword label alignment strategies, and the implementation of advanced loss functions and decoding layers to address class imbalance and sequence consistency.

\subsection{Foundational Transformer Models}
To tackle the NER task, we leverage the power of the Transformer architecture by fine-tuning three pre-trained language models. Unlike traditional methods, these models utilize the self-attention mechanism to capture deep bidirectional context:

\begin{itemize}
	\item \textbf{BERT (Bidirectional Encoder Representations from Transformers):} BERT marked a paradigm shift in NLP by capturing deep bidirectional context from unlabelled text.
	\begin{itemize}
		\item \textbf{Original Publication:} \cite{devlin2019bert}.
		\item \textbf{Architecture:} An Encoder-only Transformer. Standard configurations include $\text{BERT}_{\text{BASE}}$ (12 layers, 768 hidden units, 12 attention heads, $\approx 110\text{M}$ parameters).
		\item \textbf{Tokenization \& Data:} Utilizes the WordPiece tokenizer ($\approx 30,000$ tokens). Pre-trained on a $16\text{GB}$ dataset consisting of the Toronto BookCorpus and English Wikipedia.
		\item \textbf{Core Pre-training Strategies:} 
		\begin{itemize}
			\item \textit{Masked Language Modeling (MLM):} $15\%$ of input tokens are randomly masked. The model predicts these tokens based on context from both directions using Static Masking.
			\item \textit{Next Sentence Prediction (NSP):} A binarized task predicting whether sentence $B$ logically follows sentence $A$, capturing sentence-level cohesion: $P(\text{IsNext} | A, B) \in \{0, 1\}$.
		\end{itemize}
	\end{itemize}
	
	\vspace{0.2cm}
	\item \textbf{DistilBERT (A Distilled Version of BERT):} DistilBERT demonstrates that the size of a large language model can be significantly compressed to accelerate processing speed without severely compromising its contextual understanding capabilities.
	\begin{itemize}
		\item \textbf{Original Publication:} \cite{sanh2019distilbert}.
		\item \textbf{Optimized Compression:} DistilBERT retains the general architecture of BERT but removes token-type embeddings and the pooler, while reducing the number of encoder layers from 12 to 6. Consequently, the model is $40\%$ smaller and offers $60\%$ faster inference speed, yet retains approximately $97\%$ of the original model's performance on standard language understanding benchmarks.
		\item \textbf{Core Training Strategy (Knowledge Distillation):} Rather than proposing a novel architecture, DistilBERT optimizes efficiency through knowledge distillation, where a compact "student" model is trained to reproduce the behavior of a larger "teacher" model (BERT).
		\begin{itemize}
			\item \textit{Soft Targets:} Instead of training exclusively on strict data labels (hard targets), the student model is forced to mimic the soft probability predictions generated by the teacher. By leveraging the teacher's full probability distribution, the student acquires robust generalization capabilities without requiring a massive parameter count.
			\item \textit{Triple Loss Function:} To maximize the efficacy of knowledge transfer, DistilBERT is optimized via a linear combination of three loss components:
			\begin{enumerate}
				\item \textbf{Masked Language Modeling (MLM) Loss:} Evaluates the student's accuracy in predicting masked tokens within the text.
				\item \textbf{Distillation Loss:} Employs cross-entropy to measure and minimize the divergence between the probability distributions of the student and the teacher, compelling the student to accurately simulate the output probabilities.
				\item \textbf{Cosine-Distance Loss:} Computes the cosine distance to align the directional vectors of the hidden states between the student and the teacher, ensuring that the internal feature representations maintain a geometric structure analogous to the original model.
			\end{enumerate}
		\end{itemize}
	\end{itemize}
	
	\vspace{0.2cm}
	\item \textbf{RoBERTa (Robustly Optimized BERT Pretraining Approach):} RoBERTa demonstrates that the original BERT architecture was significantly undertrained and can achieve state-of-the-art performance simply by refining the pre-training methodology.
	\begin{itemize}
		\item \textbf{Original Publication:} \cite{liu2019roberta}.
		\item \textbf{Massive Training Corpora:} Utilizes over $160\text{GB}$ of uncompressed text (approximately $10\times$ more than BERT), incorporating massive datasets such as CC-News, OpenWebText, and Stories.
		\item \textbf{Tokenization:} Replaces BERT's WordPiece with a Byte-Level Byte Pair Encoding (BPE) algorithm. This expands the vocabulary size to roughly $50,000$ tokens, enhancing the handling of rare and out-of-vocabulary words.
		\item \textbf{Core Optimization Strategies (Refinement over Revolution):}
		\begin{itemize}
			\item \textit{Removal of Next Sentence Prediction (NSP):} Empirical results indicated that the NSP objective did not significantly improve downstream task performance. RoBERTa eliminates NSP to focus computational resources entirely on Masked Language Modeling.
			\item \textit{Dynamic Masking:} Instead of the static masking applied once during data preprocessing in BERT, RoBERTa dynamically generates masking patterns for each epoch. This stochasticity prevents the model from memorizing fixed patterns and improves generalization.
			\item \textit{High-Intensity Training:} The model is trained for more steps with significantly larger mini-batch sizes, leading to better optimization and convergence.
		\end{itemize}
	\end{itemize}
\end{itemize}

\subsection{Tokenization and Label Alignment}
Transformer models utilize subword tokenization algorithms (e.g., WordPiece or BPE). This introduces a misalignment: a single word may be split into multiple subwords (e.g., "Hoang" $\rightarrow$ ["Hoa", "\#\#ng"]). We implement a strict label alignment strategy:
\begin{enumerate}
    \item \textbf{Extraction of \texttt{word\_ids}:} Mapping each subword back to its original word index.
    \item \textbf{First-Subword Labeling:} Only the first subword of a fragmented word is assigned the original NER label (e.g., $B-SALARY$).
    \item \textbf{Ignore Index (-100):} Subsequent subwords and special tokens (e.g., $[CLS]$, $[SEP]$) are assigned a label value of $-100$, which is ignored during loss calculation in PyTorch.
\end{enumerate}

\subsection{Loss Functions and Sequence Optimization}
To mitigate class imbalance, enforce structural consistency, and handle subword tokenization artifacts (by ignoring index $-100$), we evaluate the following loss functions and decoding algorithms:

\subsubsection{Cross-Entropy Loss (CE)}
The standard distribution-based loss measuring the difference between predicted and actual distributions for valid tokens:
\begin{equation}
	L_{CE} = -\sum_{i=1}^{C} y_i \log(p_i)
\end{equation}
Where $C$ is the number of classes, $y_i$ is the binary indicator, and $p_i$ is the predicted probability.

\subsubsection{Focal Loss}
An enhancement of CE designed to handle extreme class imbalance (e.g., the dominance of the 'O' tag). It introduces a modulating factor $(1 - p_t)^\gamma$ to down-weight easy-to-classify examples, alongside a class-balancing weight $\alpha_t$ computed from the training set distribution:
\begin{equation}
	L_{Focal} = -\alpha_t (1 - p_t)^\gamma \log(p_t)
\end{equation}

\subsubsection{Dice Loss}
A region-based optimization, adapted into a continuous "Soft Dice" formulation, that directly maximizes the overlap between the predicted probabilities and the one-hot encoded ground truth:
\begin{equation}
	L_{Dice} = 1 - \frac{2 \sum p_i y_i + \sigma}{\sum p_i + \sum y_i + \sigma}
\end{equation}
Where $\sigma$ is a smoothing factor (set to $1e-6$) to prevent division by zero and stabilize gradients.

\subsubsection{Conditional Random Fields (CRF)}
Unlike independent token predictions, the CRF layer models global sequence dependencies. It learns transition probabilities to ensure output sequences follow logical naming conventions (e.g., preventing an I-tag from appearing without a preceding B-tag). The CRF layer minimizes the Negative Log-Likelihood (NLL):
\begin{equation}
	L_{CRF} = -\log P(\mathbf{y}|\mathbf{h}) = -\log \left( \frac{\exp(\text{Score}(\mathbf{h}, \mathbf{y}))}{\sum_{\mathbf{y'} \in \mathbf{Y}} \exp(\text{Score}(\mathbf{h}, \mathbf{y'}))} \right)
\end{equation}

\subsubsection{Hybrid Approaches}
To leverage the complementary strengths of the aforementioned optimization strategies, we implemented two combined architectures:
\begin{itemize}
    \item \textbf{CE + Dice Loss:} Combines the distributional stability of standard Cross-Entropy with the boundary-matching focus of Dice Loss. In our implementation, the total loss is computed as the direct sum of both components:
    $$L_{Total} = L_{CE} + L_{Dice}$$
    
    \item \textbf{CRF + Focal Loss:} Integrates the global sequence constraints of CRF with the hard-example mining capability of Focal Loss. Instead of applying Focal Loss at the token level, this hybrid calculates the sequence-level Negative Log-Likelihood (NLL) from the CRF layer and applies a focal modulating factor. Let $p_t = \exp(-\text{NLL})$ be the probability of the true label sequence, the loss is formulated as:
    $$L_{CRF-Focal} = \frac{1}{N} \sum \left( (1 - p_t)^\gamma \cdot \text{NLL} \right)$$
\end{itemize}
